\section{Hermite--Pad\'e tails}\label{sec:hp}
This section proves the two tails from the displayed construction, without
using an irrationality theorem as input.  The base is fixed throughout each
limit; constants in contour and remainder estimates may depend on that base.

\begin{theorem}\label{thm:hp-tail}
For every fixed integer $a\ge12$ or $a\le-8$, the determinant construction
below proves that $\operatorname{Li}_2(1/a)$ is irrational.
\end{theorem}

\subsection{Residues, tails, and exact integerization}
For $n\ge0$ put
\begin{equation}\label{eq:hp-phi}
 \Phi_n(w)=\frac{\prod_{r=n+1}^{3n+1}(w-r)}
                   {\prod_{j=0}^{n}(w-j)^2}
 =\sum_{j=0}^{n}\left(\frac{u_{n,j}}{(w-j)^2}
                         +\frac{v_{n,j}}{w-j}\right).
\end{equation}
Writing $H_k=\sum_{r=1}^k1/r$ and $H_0=0$, direct residue calculation gives
\begin{equation}\label{eq:hp-residues}
 u_{n,j}=-\frac{(3n+1-j)!}{(j!)^2((n-j)!)^3},\qquad
 v_{n,j}=u_{n,j}\bigl(3H_{n-j}-2H_j-H_{3n+1-j}\bigr).
\end{equation}
Set
\[
 B_n(z)=\sum_{j=0}^n v_{n,j}z^j,\qquad
 C_n(z)=\sum_{j=0}^n u_{n,j}z^j,
\]
and
\[
 A_n(z)=-\sum_{k=1}^{n}z^k\sum_{j=0}^{k-1}
       \left(\frac{v_{n,j}}{k-j}+\frac{u_{n,j}}{(k-j)^2}\right).
\]
The coefficient of $z^k$ in $B_n(-\log(1-z))+C_n\operatorname{Li}_2(z)$ is
the partial-fraction sum $\sum_{j<k}(v_{n,j}/(k-j)+u_{n,j}/(k-j)^2)$.
For $1\le k\le n$ it is canceled by $A_n$; for $n<k\le3n+1$ it is zero
because $\Phi_n(k)=0$; for $k\ge3n+2$ it equals $\Phi_n(k)$.  Therefore
\begin{equation}\label{eq:hp-tail-series}
 R_n(z):=A_n(z)+B_n(z)(-\log(1-z))+C_n(z)\operatorname{Li}_2(z)
       =\sum_{k\ge3n+2}\Phi_n(k)z^k.
\end{equation}
The adjacent Casoratian
\begin{equation}\label{eq:hp-casoratian}
 E_n(z):=B_{n+1}(z)R_n(z)-B_n(z)R_{n+1}(z)
       =W_n(z)+V_n(z)\operatorname{Li}_2(z)
\end{equation}
has
\[
 W_n=B_{n+1}A_n-B_nA_{n+1},\qquad
 V_n=B_{n+1}C_n-B_nC_{n+1}.
\]
The logarithm cancels identically.  Both $W_n,V_n$ have degree at most
$N=2n+1$.

Here is an all-degree proof of the coefficient formula needed for
integerization.  Let $b_n=\binom{2n+1}{n}$ and
\[
 J_n(t)=\sum_{k=0}^{N}(-1)^k\binom Nk\binom{n+k}{k}^{\!2}t^k.
\]
For $\Re(s)>-1$ the integral is absolutely convergent, and direct
termwise integration gives
\begin{equation}\label{eq:hp-mellin}
 \int_0^1t^sJ_n(t)\,dt=
 \frac{(2n+1)!}{(n!)^2}
 \frac{[s(s-1)\cdots(s-n+1)]^2}
      {(s+1)(s+2)\cdots(s+2n+2)}.
\end{equation}
The right side then supplies the meromorphic continuation in $s$; no
continuation is used before the stated half-plane of convergence.
Write $V_n(z)=\sum_{i=0}^Nv_i z^i$ and
$\Pi_n(t)=t^NV_n(1/t)=\sum_{j=0}^N\pi_jt^j$, so $\pi_j=v_{N-j}$.
The polynomial part of $E_n$ has degree at most $N$.  For
$k=N+\ell$, $1\le\ell\le n$, its coefficient is therefore
\[
 0=\sum_{i=0}^N\frac{v_i}{(k-i)^2}
   =\sum_{j=0}^N\frac{\pi_j}{(\ell+j)^2},
\]
which is $-\int_0^1t^{\ell-1}\Pi_n(t)\log t\,dt=0$.  The second
elimination, $C_{n+1}R_n-C_nR_{n+1}$, cancels $\operatorname{Li}_2$;
its nonpolynomial part is $V_n(z)\log(1-z)$.  The same coefficient range
gives
\[
 0=\sum_{i=0}^N\frac{v_i}{k-i}
   =\sum_{j=0}^N\frac{\pi_j}{\ell+j}
   =\int_0^1t^{\ell-1}\Pi_n(t)\,dt.
\]
Thus the two exact moment families are
\begin{equation}\label{eq:hp-moments}
 \int_0^1t^j\Pi_n(t)\,dt=0,\qquad
 \int_0^1t^j\Pi_n(t)\log t\,dt=0\quad(0\le j<n).
\end{equation}
The Mellin transform of any
polynomial of degree at most $N$ has denominator
$\prod_{r=1}^{N+1}(s+r)$ and numerator degree at most $N$.  The first family
in \eqref{eq:hp-moments} makes this numerator vanish at $s=0,\ldots,n-1$;
the second makes each zero double.  Hence the numerator is divisible by
$[s(s-1)\cdots(s-n+1)]^2$, leaving a factor of degree at most one.  The
solution space of \eqref{eq:hp-moments} consequently has dimension at most
two.

Equation~\eqref{eq:hp-mellin} shows that $J_n$ is one solution.  It also
shows $J_n(1)=0$, since the right side is $O(s^{-2})$ as $s\to+\infty$.
Integration by parts then shows that $tJ_n'(t)$ is a second solution: the
boundary terms vanish at both endpoints, and the two moment families for
$tJ'_n$ reduce to those for $J_n$.  The two solutions are independent by
their top two coefficients: if $a_N=[t^N]J_n$ and
$a_{N-1}=[t^{N-1}]J_n$, the determinant formed from the top two coefficients
of $J_n$ and $tJ_n'$ is
$a_N(N-1)a_{N-1}-Na_Na_{N-1}=-a_Na_{N-1}\ne0$.  Thus every solution is
their linear combination.  Direct substitution of the residues
\eqref{eq:hp-residues}
gives the constant and linear coefficients of $\Pi_n$; matching those two
coefficients in this two-dimensional space yields
\begin{equation}\label{eq:hp-vpoly}
 V_n(z)=b_n^2z^NP_n(1/z),
\end{equation}
where
\begin{equation}\label{eq:hp-Pn}
 P_n(t)=-\sum_{k=0}^{N}(-1)^k\binom Nk\binom{n+k}{k}
 \left((12n+17)\binom{n+k}{k}
       +(21n+29)\binom{n+k}{k-1}\right)t^k.
\end{equation}
The convention is $\binom{n+k}{k-1}=0$ at $k=0$.  In particular, the
second binomial is $\binom{n+k}{k-1}$, not
$\binom{n+k-1}{k-1}$.

The first two coefficients follow from the endpoint residues. Put
$\eta_m=-2H_m-H_{2m+1}$, $\rho_m=2m^2(m+1)$, and
$\kappa_m=3+2/m-1/(2m+2)$.  Their ratios are
\[
 \frac{v_{m,m}}{u_{m,m}}=\eta_m,\qquad
 \frac{u_{m,m-1}}{u_{m,m}}=\rho_m,\qquad
 \frac{v_{m,m-1}}{u_{m,m-1}}=\eta_m+\kappa_m.
\]
For $n\ge1$, expanding $V_n=B_{n+1}C_n-B_nC_{n+1}$ at its top two
degrees yields, with $\pi_j=[t^j]\Pi_n$,
\begin{align*}
 \pi_0&=u_{n,n}u_{n+1,n+1}(\eta_{n+1}-\eta_n),\\
 \pi_1&=u_{n,n}u_{n+1,n+1}
 \bigl((\rho_n+\rho_{n+1})(\eta_{n+1}-\eta_n)
       -\rho_n\kappa_n+\rho_{n+1}\kappa_{n+1}\bigr).
\end{align*}
Here $u_{n,n}=-(n+1)b_n$, $u_{n+1,n+1}=-2(2n+3)b_n$, and
$\eta_{n+1}-\eta_n=-5/(2(n+1))-1/(2n+3)$.  Substitution yields
$\pi_0=-b_n^2(12n+17)$ and
\[
 \pi_1=b_n^2(2n+1)(n+1)
       \bigl((n+1)(12n+17)+(21n+29)\bigr).
\]
For $n=0$, the residues give the same formulas directly: $\pi_0=-17$ and
$\pi_1=46$.  For the Casoratian calculation below, put
$h_m=3H_m-H_{3m+1}$ and
\[
 \delta_m=h_{m+1}-h_m=
 \frac{2(27m^2+54m+23)}{3(m+1)(3m+2)(3m+4)},
\]
The values of $\pi_0,\pi_1$ are the first two coefficients of
$-(12n+17)J_n-(21n+29)tJ'_n/(n+1)$.  The identity
$\frac{k}{n+1}\binom{n+k}{k}=\binom{n+k}{k-1}$ gives
\eqref{eq:hp-Pn}.  The dimension bound proves \eqref{eq:hp-vpoly} for every
$n$, with no finite-range restriction.

For reference, the boundary residues used in this matching can be checked
directly from \eqref{eq:hp-residues}.  For $m\ge1$,
\begin{align}
 u_{m,m-1}&=-\frac{(2m+2)!}{((m-1)!)^2},\qquad
 u_{m,m}=-\frac{(2m+1)!}{(m!)^2},\\
 \frac{u_{m,m-1}}{u_{m,m}}&=(2m+2)m^2=2m^2(m+1),\\
 \frac{v_{m,m-1}}{u_{m,m-1}}-\frac{v_{m,m}}{u_{m,m}}
   &=3+\frac2m-\frac1{2m+2},\\
 \frac{u_{n,n}u_{n+1,n+1}}{b_n^2}&=(2n+2)(2n+3).
\end{align}
These identities also cover the boundary-residue normalization in the
three-form extraction below; the case $n=0$ is read directly from
\eqref{eq:hp-residues}.

Let $d_N=\operatorname{lcm}(1,\ldots,N)$ and
$Q_n(z)=V_n(z)/b_n^2\in\mathbb Z[z]$.  Since $E_n=W_n+V_n\operatorname{Li}_2$
vanishes through degree $N$, its polynomial part satisfies
\[
 \frac{W_n(z)}{b_n^2}=-[Q_n(z)\operatorname{Li}_2(z)]_{\le N}.
\]
Each coefficient on the right has denominators among $r^2$ with
$1\le r\le N$, and $r\mid d_N$.  Therefore
\begin{equation}\label{eq:hp-integer-poly}
 \frac{d_N^2}{b_n^2}W_n,\quad
 \frac{d_N^2}{b_n^2}V_n\in\mathbb Z[z].
\end{equation}
For every fixed integer $a$ with $|a|\ge2$, put
\begin{equation}\label{eq:hp-integer-form}
 L_n(a)=a^N\frac{d_N^2}{b_n^2}E_n(1/a)
       =X_n(a)+Y_n(a)\operatorname{Li}_2(1/a),
       \qquad X_n(a),Y_n(a)\in\mathbb Z.
\end{equation}
The factor $a^N$ clears evaluation denominators because both polynomials in
\eqref{eq:hp-integer-poly} have degree at most $N$.

\subsection{Remainder bound for both signs}
For $-1<z<1$, beta integration of \eqref{eq:hp-tail-series} gives the exact
identity
\begin{equation}\label{eq:hp-beta}
 R_n(z)=\frac{(2n+1)!}{(n!)^2}z^{3n+2}
 \int_0^1\!\int_0^1\frac{tu}{(1-ztu)^2}
 \left(\frac{t^2u^2(1-t)(1-u)}{(1-ztu)^2}\right)^n dt\,du.
\end{equation}
To maximize the bracket, fix $p=tu$.  Since $t+u\ge2\sqrt p$,
$(1-t)(1-u)=1-(t+u)+p\le(1-\sqrt p)^2$, with equality only at
$t=u=\sqrt p$.  Write $s=\sqrt p$.  The bracket is at most
\[
 g_z(s)=\frac{s^4(1-s)^2}{(1-zs^2)^2},\qquad 0\le s\le1.
\]
For $|z|<1$, its unique interior maximum satisfies
$zs^3-3s+2=0$: differentiating $\log g_z$ gives
$2-3s+zs^3=0$, and the left side is strictly decreasing on $[0,1]$.
At that root, $1-zs^2=2(1-s)/s$, so $g_z(s)=s^6/4$.  The prefactor
$tu/(1-ztu)^2$ in \eqref{eq:hp-beta} is bounded for fixed $z$ with
$|z|<1$.  Stirling's formula~\cite{WhittakerWatson} and the unit area of the square therefore give
\begin{equation}\label{eq:hp-R-rate}
 \limsup_{n\to\infty}|R_n(z)|^{1/n}\le |z|^3s(z)^6.
\end{equation}
No cancellation estimate is used.

\subsection{Positive coefficient bound}
For $z=1/q$, $q>1$, all $u_{n,j}$ have the same (negative) sign.  Put
$x=j/n$ and
\[
 f_q(x)=(3-x)\log(3-x)-2x\log x-3(1-x)\log(1-x)-x\log q,
 \quad 0\le x\le1,
\]
where $0\log0=0$.  Uniform Stirling estimates~\cite{WhittakerWatson} in
\eqref{eq:hp-residues} give
\[
 \log\bigl(|u_{n,j}|q^{-j}\bigr)=nf_q(j/n)+O(\log(n+1))
\]
uniformly for $0\le j\le n$.  Indeed, the coefficients of $n\log n$ and
$n$ cancel, and all factorial contributions are $O(\log(n+1))$,
including at endpoints.  Now
\[
 f_q'(x)=\log\frac{(1-x)^3}{q x^2(3-x)},\qquad
 f_q''(x)=\frac1{3-x}-\frac2x-\frac3{1-x}<0.
\]
There is a unique maximizer $t(q)\in(0,1)$, determined by
$(1-t)^3=qt^2(3-t)$.  Substitution of this equation gives
\[
 \exp(f_q(t))=\left(\frac{3-t}{1-t}\right)^3.
\]
Since the sum has only $n+1$ terms,
\[
 \limsup_{n\to\infty}|C_n(1/q)|^{1/n}
 \le \left(\frac{3-t(q)}{1-t(q)}\right)^3.
\]
Also $|3H_{n-j}-2H_j-H_{3n+1-j}|\le6(1+\log(3n+1))$, so the same bound
holds for $B_n(1/q)$.  This upper-bound form of the discrete Laplace
principle is all that is needed; no unproved ratio limit is used.

\subsection{Negative coefficient bound: contour and endpoint estimate}
Fix $q>1$, put $z=-1/q$, and define the two logarithmic branches
\[
 F_n^\pm(z)=B_n(z)+C_n(z)(-\log q\pm i\pi).
\]
Residues of $\Phi_n(w)e^{(-\log q\pm i\pi)w}$ give
\begin{equation}\label{eq:hp-contour-res}
 F_n^\pm(z)=\frac1{2\pi i}\oint_{\Gamma_n^\pm}
          \Phi_n(w)e^{(-\log q\pm i\pi)w}\,dw,
\end{equation}
where the rectangle encloses the poles $0,1,\ldots,n$.  Set $\beta_q>0$
equal to the absolute imaginary part of the lower nonreal root of
\begin{equation}\label{eq:hp-saddle-cubic}
 p_q(w)=(q+1)w^3-3(q+1)w^2+3w-1=0.
\end{equation}
For all sufficiently large $n$, let $\Gamma_n^+$ be the positively oriented
rectangle with vertical sides $\Re w=-1/2,n+1/2$ and horizontal sides
$\Im w=\pm n\beta_q$.  It encloses precisely the poles $0,1,\ldots,n$.
Its distance from each enclosed pole is at least
$\min\{1/2,n\beta_q\}$, hence at least $1/2$ once
$n\beta_q\ge1/2$.  All numerator zeros lie outside the rectangle at distance
at least $1/2$.  The rate proof uses only these sufficiently large indices.
For $F_n^-$ use the conjugate exponential branch on the same positively
oriented rectangle.

We give the uniform sum-to-integral estimate needed on these contours.
If $\xi=w/n$ lies on the scaled rectangle, then
\begin{equation}\label{eq:hp-log-em}
 \log|\Phi_n(n\xi)|=n\left(
 \int_1^3\log|\xi-r|\,dr-2\int_0^1\log|\xi-r|\,dr\right)
 +O_q(\log(n+1)),
\end{equation}
uniformly along all four sides.  On a horizontal side $|\Im\xi|=\beta_q$,
the functions $r\mapsto\log|\xi-r|$ have uniformly bounded derivative, so
the ordinary left/right Riemann-sum error is $O_q(1)$.  On the left
vertical side $\Re\xi=-1/(2n)$, the real part lies outside both integration
intervals by at least $1/(2n)$.  On the right side
$\Re\xi=1+1/(2n)$ it lies outside $[0,1]$ by the same amount.  For $[1,3]$
there, the only possible near-singularity lies in the first mesh cell:
its real part is exactly halfway between the endpoint and the first mesh
point.  If $|\Im\xi|\le1/n$, rescaling this one cell by $r=1+u/n$ shows
that both its mesh contribution and $n$ times its integral are
$O(\log(n+1))$; the integral of $\log\sqrt{u^2+v^2}$ over a fixed
interval is bounded uniformly for $0\le v\le1$.  If
$|\Im\xi|>1/n$, or on the remaining cells, the total variation is bounded
by $C_q+\int_{1/n}^{O_q(1)}dr/r=O_q(\log(n+1))$.  Endpoint shifts of one
mesh point in the numerator product contribute at most another
$O_q(\log(n+1))$.  This proves \eqref{eq:hp-log-em}; in particular the
logarithmic singularity is handled explicitly rather than hidden in a
smooth-function Euler--Maclaurin assertion.

For completeness, the estimate in \eqref{eq:hp-log-em} is a bounded-variation
Riemann-sum estimate, including the endpoint cell.  Put
$f_\xi(r)=\log|\xi-r|$.  The numerator sum is the right Riemann sum on
$[1,3]$, with the extra point $f_\xi(3+1/n)$, and the denominator sum is the
right Riemann sum on $[0,1]$, with $f_\xi(0)$.  On a horizontal edge the
total variation is $O_q(1)$.  On a vertical edge the only possible singular
cell is the first half-mesh cell; after $r=1+u/n$ (or $r=u/n$) its mesh term
and its $n$-scaled integral are $O(\log(n+1))$.  The remaining variation is
$O_q(1)+\int_{1/n}^{O_q(1)}dr/r=O_q(\log(n+1))$.  Thus the estimate is
uniform on all four edges and does not invoke a smooth Euler--Maclaurin
formula across a logarithmic endpoint.

Combining \eqref{eq:hp-log-em} with the exact modulus of the exponential
factor gives
\[
 \log|\Phi_n(n\xi)e^{(-\log q+i\pi)n\xi}|
       =nH_q(\xi)+O_q(\log(n+1)),
\]
where
\begin{equation}\label{eq:hp-H}
 H_q(\xi)=\int_1^3\log|\xi-r|\,dr-2\int_0^1\log|\xi-r|\,dr
          -\log q\,\Re\xi-\pi\Im\xi.
\end{equation}
The discriminant of \eqref{eq:hp-saddle-cubic} is
$-108q^2(q+1)<0$, and
$p_q(2)=1-4q<0<p_q(3)=8$.  Thus the real root is in $(2,3)$; the
conjugate roots have real part $u=(3-R)/2\in(0,1/2)$.  Write the lower
root as $\xi_0=u-i\beta_q$.  Differentiating the analytic logarithmic
phase underlying \eqref{eq:hp-H} shows $\partial_xH_q(\xi_0)=0$ by the
saddle equation $z(\xi-1)^3=\xi^2(\xi-3)$.  Along the lower side,
\[
 \frac{d^2}{dx^2}H_q(x-i\beta_q)
 =-6\Re\frac1{(x-i\beta_q)(x-1-i\beta_q)(x-3-i\beta_q)}<0
 \quad(0\le x\le1),
\]
because the real part of the denominator is
$x(x-1)(x-3)+\beta_q^2(4-3x)>0$.  Hence the unique lower-side maximum is
$\xi_0$.  On the vertical sides, parameterized downward,
\[
 \frac d{dy}H_q(-iy)=3\arctan y-\arctan(y/3)>0,
\]
\[
 \frac d{dy}H_q(1-iy)=\frac\pi2+2\arctan y-\arctan(y/2)>0.
\]
On the lower halves of the limiting vertical sides, parameterized as
$\xi=-iy$ and $\xi=1-iy$ with $0<y\le\beta_q$, the displayed derivatives
are strictly positive.  On the upper halves use the exact reflection identity
$H_q(x+iy)=H_q(x-iy)-2\pi y$ for $y>0$.  Hence neither upper vertical half
can exceed its reflected lower half, and the maxima on both lower vertical
halves occur at their lower endpoints.  The strictly concave lower horizontal
phase has its unique maximum at the interior saddle $\xi_0$, strictly above
both endpoints.

The logarithmic potentials defining $H_q$ are continuous when $\xi$ approaches
the real integration intervals because their logarithmic singularities are
integrable.  They are therefore uniformly continuous on a compact
neighborhood of the limiting rectangle.  The maxima over the boundaries of
the scaled finite rectangles converge to the maximum over the limiting
boundary.  The original rectangles have perimeter $O_q(n)$, while
\eqref{eq:hp-log-em} contributes only a polynomial factor in $n$ to the
modulus bound.  Taking $n$-th roots now gives
the claimed rate bound.
\begin{equation}\label{eq:hp-F-rate}
 \limsup_{n\to\infty}|F_n^\pm(-1/q)|^{1/n}\le e^{H_q(\xi_0)}.
\end{equation}

To evaluate this number, choose logarithms continuously from the lower
contour.  An antiderivative gives the analytic phase
\[
 \Psi_q(\xi)=3(\xi-1)\operatorname{Log}(\xi-1)- (\xi-3)\operatorname{Log}(\xi-3)
             -2\xi\operatorname{Log}\xi+(-\log q+i\pi)\xi.
\]
At the lower saddle choose the principal logarithms of
$\xi_0,\xi_0-1,\xi_0-3$, whose arguments lie in $(-\pi,0)$ (with the
negative-real limiting value approached from below).  The algebraic saddle
equation implies that $\Psi_q'(\xi_0)$ is an integer multiple of $2\pi i$.
As $q\to\infty$, the lower root satisfies
$\xi_0=-i/\sqrt{3q}+O(q^{-1})$; the three arguments tend respectively to
$-\pi/2,-\pi,-\pi$.  The corresponding imaginary-part combination is zero,
so continuity in $q>1$ forces this integer multiple to vanish identically.
Substituting the saddle equation into $\Psi_q$ gives
\[
 e^{H_q(\xi_0)}=\left|\frac{\xi_0-3}{\xi_0-1}\right|^3.
\]
For $y=(\xi-3)/(\xi-1)$ the saddle equation becomes
$y(y-3)^2=4z$.  Its real root is $y_0=-s_-^2/q$, where
$s_-(q)\in(0,2/3)$ is the unique solution of
\begin{equation}\label{eq:hp-sminus}
 -\frac{s_-^3}{q}-3s_-+2=0.
\end{equation}
The product of the three roots of $y(y-3)^2=4z$ is $4z=-4/q$.
Consequently the squared modulus of either nonreal root is
$4/s_-^2$, and
\begin{equation}\label{eq:hp-negative-coeff-rate}
 e^{H_q(\xi_0)}=\frac8{s_-(q)^3}.
\end{equation}
Finally
\[
 C_n(-1/q)=\frac{F_n^+-F_n^-}{2\pi i},\qquad
 B_n(-1/q)=\frac{F_n^++F_n^-}{2}+\log q\,C_n(-1/q).
\]
Equations~\eqref{eq:hp-F-rate}--\eqref{eq:hp-negative-coeff-rate} thus give
the same upper exponential rate $8/s_-^3$ for $B_n(-1/q)$.

\subsection{Integer-form rates and exact cutoffs}
The prime number theorem gives
$\log d_N=N+o(N)$, while Stirling~\cite{WhittakerWatson} gives
$n^{-1}\log b_n\to2\log2$.  With $N=2n+1$,
\begin{equation}\label{eq:hp-den-rate}
 \lim_{n\to\infty}\left(\frac{d_N^2}{b_n^2}\right)^{1/n}
       =\frac{e^4}{16}.
\end{equation}
No finite global inequality for $d_N$ is asserted.  Combining
\eqref{eq:hp-R-rate}, the coefficient bounds, and
\eqref{eq:hp-casoratian} yields
\begin{align}
 \limsup |L_n(q)|^{1/n}&\le
 r_+(q):=\frac{e^4}{16q}
       \left(\frac{3-t(q)}{1-t(q)}\right)^3s_+(q)^6,
       &&a=q>1,\\
 \limsup |L_n(-q)|^{1/n}&\le
 r_-(q):=\frac{e^4s_-(q)^3}{2q},
       &&a=-q<-1,
\end{align}
where $s_+$ is the unique root in $(0,1)$ of
$s_+^3/q-3s_++2=0$.  The factors are as follows: $q^2$ from $|a|^N$,
$e^4/16$ from \eqref{eq:hp-den-rate}, $q^{-3}s^6$ from the remainder,
and the relevant coefficient growth rate.  Thus no estimate for an
uneliminated logarithmic form is being used.

For the positive rate set $z=1/q$.  Implicit differentiation of the two
root equations gives
\[
 \frac{d\log r_+}{d\log z}
       =1+t+\frac{2zs_+^2}{1-zs_+^2}>0,
\]
so $r_+(q)$ decreases with $q$.  On the negative side,
$q=s_-^3/(2-3s_-)$ is increasing in $s_-$ and
$r_-(q)=e^4(2-3s_-)/2$ is decreasing.  It remains to check four rational
endpoints.  The Taylor sum for $e^4$ through degree 16 and a geometric tail
bound give
\[
 \frac{6972316451}{127702575}<e^4<\frac{377139364267}{6907548375}.
\]
Opposite signs of the defining cubics at the following rational endpoints
isolate the roots.  The derivatives are
$-3(1-t)^2-q(6t-3t^2)<0$ for $t(q)$,
$3s^2/q-3<0$ for $s_+(q)$, and $-3s^2/q-3<0$ for $s_-(q)$:
\[
\begin{array}{c|c}
 t(11)&(14179/100000,14180/100000)\\
 s_+(11)&(67602/100000,67604/100000)\\
 t(12)&(13681/100000,13683/100000)\\
 s_+(12)&(67521/100000,67523/100000)\\
 s_-(7)&(65338/100000,65339/100000)\\
 s_-(8)&(65496/100000,65497/100000)
\end{array}
\]
Exact rational interval substitution gives
\[
\begin{array}{c|c}
 r_+(11)&(10937/10000,10941/10000)\\
 r_+(12)&(9834/10000,9837/10000)\\
 r_-(7)&(10877/10000,10879/10000)\\
 r_-(8)&(9587/10000,9589/10000)
\end{array}
\]
and hence $r_+(12)<1$ and $r_-(8)<1$.  Monotonicity proves
$r_+(q)<1$ for every integer $q\ge12$ and $r_-(q)<1$ for every integer
$q\ge8$.

\subsection{All-degree nonvanishing and conclusion}
It remains to show that the constructed integer forms do not vanish
eventually.
Let $U_m=(3m+1)!/(m!)^3$,
$K_n=((2n+1)!)^3/((3n+2)!)^2$, and
$h_m=3H_m-H_{3m+1}$.  Direct extraction from
\eqref{eq:hp-tail-series} gives
\[
 \Phi_n(3n+2)=K_n,\qquad
 \frac{\Phi_n(3n+3)}{K_n}=\frac{8(n+1)}9.
\]
For $m\ge0$, put $M_m=B_mC_{m+1}-C_mB_{m+1}$; in particular
$M_0=B_0C_1-C_0B_1$ is defined by the displayed residue polynomials.
For $m\ge1$, put $r_m=m^3/(3m+1)$ and $c_m=2+3/m-1/(3m+1)$.
The lower endpoint residues give
\[
 u_{m,0}=-U_m,\quad u_{m,1}=r_m u_{m,0},\quad
 v_{m,0}=h_m u_{m,0},\quad
 v_{m,1}=r_m(h_m-c_m)u_{m,0}.
\]
Expanding $B_mC_{m+1}-C_mB_{m+1}$ at degrees zero and one gives
\[
 M_m(0)=-U_mU_{m+1}\delta_m,
\]
\[
 [z]M_m=U_mU_{m+1}
 \{-\delta_m(r_m+r_{m+1})+r_{m+1}c_{m+1}-r_mc_m\},
\]
The three-form determinant
\[
 \Delta_n(z)=\det\begin{bmatrix}
 A_n&B_n&C_n\\ A_{n+1}&B_{n+1}&C_{n+1}\\ A_{n+2}&B_{n+2}&C_{n+2}
 \end{bmatrix}
\]
has degree at most $3n+3$.  Replacing its first column by the corresponding
$R$-column and using \eqref{eq:hp-tail-series} shows its order is at least
$3n+2$.  Expanding along that column gives
\[
 \Delta_n=R_nM_{n+1}
 -R_{n+1}(B_nC_{n+2}-C_nB_{n+2})+R_{n+2}M_n.
\]
The last two terms start at degree $3n+5$.  Also
\[
 R_n(z)=K_nz^{3n+2}\left(1+\frac{8(n+1)}9z+O(z^2)\right).
\]
It follows that
\[
 [z^{3n+2}]\Delta_n=K_nM_{n+1}(0),\qquad
 \frac{[z^{3n+3}]\Delta_n}{[z^{3n+2}]\Delta_n}
 =\frac{[z]M_{n+1}}{M_{n+1}(0)}+\frac{8(n+1)}9.
\]
Substitution of the displayed coefficients gives the ratio $-\theta_n$, with
\[
 \theta_n=
 \frac{-\delta_{n+1}(r_{n+1}+r_{n+2})
       +r_{n+2}c_{n+2}-r_{n+1}c_{n+1}}{\delta_{n+1}}
       -\frac{8(n+1)}9.
\]
Substituting $r_m,c_m,\delta_m$ and clearing denominators gives
\begin{equation}\label{eq:hp-theta}
 \theta_n=\frac{1728n^4+11520n^3+28140n^2+29965n+11770}
 {18(3n+4)^2(27n^2+108n+104)}.
\end{equation}
Therefore
\begin{equation}\label{eq:hp-delta}
 \Delta_n(z)=-K_nU_{n+1}U_{n+2}\delta_{n+1}z^{3n+2}(1-\theta_nz),
\end{equation}
The numerator is positive for $n\ge0$; denominator minus numerator is
\[
 2646n^4+17640n^3+43140n^2+46067n+18182>0.
\]
Also $\delta_{n+1}>0$.  Therefore $\Delta_n(z)\ne0$ for every real
$0<|z|<1$.

For such $z$, the beta integral \eqref{eq:hp-beta} gives $R_m(z)\ne0$.
If both $E_n(z)$ and $E_{n+1}(z)$ vanished, the two adjacent minors of the
columns $A+\operatorname{Li}_2(z)C$ and $B$ would vanish.  Since
$R_n(z)$ and $R_{n+1}(z)$ are nonzero, this makes the three columns in
$\Delta_n(z)$ linearly dependent (the conclusion also holds when one of the
corresponding $B$ entries is zero, by using the other adjacent equation).
Hence $\Delta_n(z)=0$, contradicting \eqref{eq:hp-delta}.  Thus no two
consecutive $E_n(z)$ vanish.

Suppose now that $\operatorname{Li}_2(1/a)=r/s$ were rational, in lowest
terms.  By \eqref{eq:hp-integer-form}, $sL_n(a)$ is an integer for every
$n$.  The strict rate bounds imply $L_n(a)\to0$, so $sL_n(a)=0$ for all
sufficiently large $n$.  This contradicts the fact that no two consecutive
$E_n(1/a)$, and hence no two consecutive $L_n(a)$, vanish.  This proves
Theorem~\ref{thm:hp-tail}.
